%======================================================================
% REFERENCIAS
%======================================================================
\chapter{Referencias}
\addcontentsline{toc}{chapter}{Referencias}

\bigskip

\begin{thebibliography}{99}

\bigskip

bibitem{amdahl}
Amdahl, G. M. (1967).
\newblock Validity of the single processor approach to achieving large scale computing capabilities.
\newblock {\em Proceedings of the AFIPS Spring Joint Computer Conference}, 483-485.

\bigskip

bibitem{rust}
The Rust Programming Language.
\newblock {\em https://www.rust-lang.org/}

\bigskip

bibitem{tokio}
Tokio: A runtime for writing reliable asynchronous applications with Rust.
\newblock {\em https://tokio.rs/}

\bigskip

bibitem{rayon}
Rayon: A data parallelism library for Rust.
\newblock {\em https://github.com/rayon-rs/rayon}

\bigskip

bibitem{rhai}
Rhai: A lightweight and fast scripting engine for Rust.
\newblock {\em https://rhai.rs/}

\bigskip

bibitem{serde}
Serde: Serialization framework for Rust.
\newblock {\em https://serde.rs/}

\bigskip

bibitem{grpc}
gRPC: A high-performance, open-source universal RPC framework.
\newblock {\em https://grpc.io/}

\bigskip

bibitem{iso25010}
ISO/IEC 25010:2011.
\newblock Systems and software engineering -- Systems and software Quality Requirements and Evaluation (SQuaRE).

\bigskip

bibitem{iso27001}
ISO/IEC 27001:2013.
\newblock Information technology -- Security techniques -- Information security management systems.

\bigskip

bibitem{chrono}
Chrono: Date and time library for Rust.
\newblock {\em https://docs.rs/chrono/}

\bigskip

bibitem{inmemory}
In-Memory Computing: Principles and Applications.
\newblock {\em Springer, 2019.}

\end{thebibliography}

\vfill
\newpage
